\documentclass[slides]{pgnotes}

\title{Non-relational databases}

\begin{document}

\maketitle

\section*{Recommended reading}

\begin{enumerate}
\item \href{https://www.mongodb.com/nosql-explained}{What is NoSQL? by MongoDB}
\item \href{https://www.zoomdata.com/master-class/architecture-and-deployment/everything-old-new-again-sql-relational-database/}{Everything Old is New Again: Before SQL Relational Database}
\item BASE: an ACID alternative \citep{pritchett:2008:base}
\end{enumerate}

\section{Historical context}

Databases have existed in various forms almost as long as mainframe computers have existed.
For our discussion here, we will assume a database to be a distinct software component that manages its own files and is used by higher-level application programs.
(i.e. we will ignore file-based storage)

\subsection{Pre-relational}

\subsubsection{Hierarchical model}

The earliest databases were based on a hierarchical model, \autoref{fig:hierarchical-model}.
IBM Information Management System (IMS) would be a typical example that is still used today. 

\autoimage{hierarchical_model}{Hierarchical model}{hierarchical-model}

\subsubsection{Network model}

The hierarchical model was extended somewhat to become a network model, \autoref{fig:network-model}, where nodes had multiple parents.

\autoimage{network_model}{Network model}{network-model}

\subsubsection{CODASYL model}

The network model became the basis of the CODASYL model, \autoref{fig:codasyl-model}, that formed the basis of a number of commerical database products.
Almost all are obsolete by this now. 

\autoimage{codasyl_model}{CODASYL model}{codasyl-model}

\newpage
\subsection{Relational databases}

By 1980's interest primarily was in relational databases, \autoref{fig:relational-model}.

\autoimage{relational_model}{Relational model}{relational-model}

\begin{itemize}
\item Largely complied (though none completely) with \href{https://en.wikipedia.org/wiki/Codd\%27s_12_rules}{Codd's 12 rules}.
\item Majority, though not all, support Structured Query Language (SQL) as their primary DML and DDL.
\end{itemize}


\subsection{Relational model limitations}

\begin{itemize}
\item Not all data fits the relational model optimally  
\item Pre-defined schema can be problematic in some contexts
  \begin{itemize}
  \item Perception (perhaps incorrect) that relational databases are ``not agile''
  \end{itemize}
\item Horizontal scaling difficult due to ACID compliance
  \begin{itemize}
  \item Read-replica can help, but not suited for write-heavy applications.
  \end{itemize}
\end{itemize}

\section{Non-relational databases}

Non-relational databases as the name suggests eschew the relational model.
Aside from that, modern non-relational databases differ widely in their data representation and access patterns.
Some characteristics are explored here.


\subsection{Query languages}

SQL is closely tied to the components and ideas of a relational database: tables, joins, foreign keys etc.
Non-relational databases are sometimes referred to as NoSQL in light of the fact that they don't use SQL. 

Most NoSQL databases follow a similar client-server pattern to their RDBMS counterparts:
\begin{itemize}
\item Database listens at a network port
\item Database may have a textual query language. Usually DB-dependent. Not standardised like SQL.
\item Bindings / drivers / libraries are usually used to access the database from programming languages.
\end{itemize}

\subsection{CAP theorem}

CAP theorem states:
\begin{quotation}
  impossible to have both consistency and availability in a partition-tolerent distributed system.
\end{quotation}

In practice:

\begin{itemize}
  \item   We have already met the ACID consistency model.
  \begin{itemize}
  \item ACID favours consistency over availability.
  \end{itemize}
  \citet{pritchett:2008:base} proposed an alternative to ACID, called BASE:
  \begin{itemize}
  \item BASE favourrs availability over consistency.
  \end{itemize}
\end{itemize}

\subsection{BASE consistency model}

\begin{description}
\item[Basic availability:]
  basic reading and writing operations are available as much as possible (using all nodes of a database cluster). No consistency guarantees:
    \begin{itemize}
    \item writes might not persist after conflicts are reconciled
    \item read may not get the latest write
    \end{itemize}
  \item[Soft-state:] within interval after write, probability $<1$ of knowing the final state, since it may not yet have converged
  \item[Eventual consistency:] when system functioning and after sufficient interval after any given set of inputs, know what the state of the database is. Any further reads will be consistent with our expectations. 
\end{description}

Most, not all, non-relational databases tend towards a BASE consistency model. 

\section{Common non-relational databases}

\subsection{Key-Value database}

A \href{https://en.wikipedia.org/wiki/Key-value_database}{key-value database} maintains a very simple mapping of keys to values.
Example: memcached, redis

\subsection{Document databases}

Model is the document.
Works well for storing data that's easily encapsulated into a single logical object.
Example: mongodb (has a lot of facilities that mirror those of a relational DB)

\subsection{Graph databases}

Model is nodes and relationships, both having attributes.
Designed for ease of traversal among nodes.
Example: neo4j

\subsection{Others}

There are many other types of NoSQL database.

\bibliographystyle{plainnat}
\bibliography{bibliography}

\end{document}
